% ===== PRÉAMBULE CANONIQUE — à coller VERBATIM en tête de CHAQUE .tex =====
\documentclass[11pt,a4paper]{article}
\usepackage[utf8]{inputenc}\usepackage[T1]{fontenc}\usepackage{lmodern}
\usepackage[french]{babel}
\usepackage{amsmath,amssymb,mathtools}\usepackage{icomma}
\usepackage[version=4]{mhchem}          % \ce{...} : équations & formules
\usepackage{siunitx}\sisetup{locale=FR} % \qty{}{}, \si{}, \ang{}
\usepackage{chemfig}                    % \chemfig{...} : structures développées
\usepackage{xcolor}\usepackage{colortbl}
\usepackage{tikz}\usepackage{pgfplots}\pgfplotsset{compat=1.18}
\usepackage[most]{tcolorbox}\usepackage{enumitem}\usepackage{array,booktabs}
\usepackage[a4paper,margin=2.2cm]{geometry}
\usetikzlibrary{arrows.meta,calc,angles,quotes,positioning,patterns,backgrounds,shapes.geometric}
\definecolor{primary}{RGB}{222,49,99}\definecolor{primarydark}{RGB}{150,22,55}
\definecolor{defcol}{RGB}{0,114,178}\definecolor{methcol}{RGB}{52,92,148}
\definecolor{excol}{RGB}{0,135,90}\definecolor{attncol}{RGB}{200,40,55}
\tcbset{boxbase/.style={enhanced,breakable,boxrule=0.9pt,arc=2.5pt,
  left=3mm,right=3mm,top=2.3mm,bottom=2.3mm,fonttitle=\bfseries,coltitle=white}}
% --- boîtes TUTORIEL (noms EXACTS, ne pas renommer) ---
\newtcolorbox{cle}[1][]{boxbase,colback=defcol!6,colframe=defcol,
  title={\textsc{À retenir}\ifx\relax#1\relax\relax\else\textmd{ : \itshape #1}\fi}}
\newtcolorbox{methode}[1][]{boxbase,colback=methcol!7,colframe=methcol,sharp corners,
  title={\textsc{Méthode}\ifx\relax#1\relax\relax\else\textmd{ --- \itshape #1}\fi}}
\newtcolorbox{exemplebox}[1][]{boxbase,colback=excol!5,colframe=excol,
  title={\textsc{Exemple}\ifx\relax#1\relax\relax\else\textmd{ : \itshape #1}\fi}}
\newtcolorbox{piege}[1][]{boxbase,colback=attncol!6,colframe=attncol,
  title={\textsc{Piège}\ifx\relax#1\relax\relax\else\textmd{ : \itshape #1}\fi}}
\newtcolorbox{remarque}[1][]{boxbase,colback=black!4,colframe=black!45,
  title={\textsc{Remarque}\ifx\relax#1\relax\relax\else\textmd{ : \itshape #1}\fi}}
% --- boîtes EXERCICE / CORRIGÉ (noms EXACTS, auto-comptées) ---
\newtcolorbox[auto counter]{exo}[1][]{boxbase,colback=primary!4,colframe=primary,
  title={\textsc{Exercice}~\thetcbcounter\ifx\relax#1\relax\relax\else\textmd{ --- \itshape #1}\fi}}
\newtcolorbox[auto counter]{corrige}[1][]{boxbase,colback=excol!4,colframe=excol,
  title={\textsc{Corrigé}~\thetcbcounter\ifx\relax#1\relax\relax\else\textmd{ --- \itshape #1}\fi}}
\newcommand{\R}{\mathbb{R}}\newcommand{\N}{\mathbb{N}}\newcommand{\Z}{\mathbb{Z}}\newcommand{\ds}{\displaystyle}
% (un \usepackage{fancyhdr}+en-tête est possible ; il est ignoré côté web)
% =========================================================================
\begin{document}
\sisetup{per-mode=symbol}
\begin{exo}[vrai ou faux, justifié]
Indique si chaque affirmation est \textbf{vraie} ou \textbf{fausse}, et justifie.
\begin{enumerate}[label=\alph*)]
  \item «~Une mesure très fidèle est forcément juste.~»
  \item «~Plus une mesure comporte de chiffres significatifs, plus son incertitude relative est faible.~»
  \item «~L'incertitude relative s'exprime dans la même unité que la grandeur mesurée.~»
  \item «~Une balance affichant $4$ chiffres significatifs ne peut pas donner une valeur fausse.~»
\end{enumerate}
\end{exo}

\begin{exo}[analyse d'une série de titrages]
Cinq titrages du même échantillon donnent les volumes équivalents
$12,45\,;\ 12,47\,;\ 12,44\,;\ 12,46\,;\ 12,48\ \si{\milli\litre}$. La valeur de référence (vraie) est
$\qty{12,46}{\milli\litre}$.
\begin{enumerate}[label=\alph*)]
  \item La série est-elle fidèle ? juste ?
  \item Calcule la valeur moyenne.
  \item Donne l'incertitude relative de la mesure (prends $\Delta V$ au rang du dernier CS).
\end{enumerate}
\end{exo}

\begin{exo}[le piège de la pesée différentielle]
On détermine la masse d'un précipité par pesée différentielle : bécher vide $m_1=\qty{125,46}{\gram}$,
puis bécher $+$ précipité $m_2=\qty{125,78}{\gram}$ (balance au centigramme, $\Delta=\qty{0,01}{\gram}$).
\begin{enumerate}[label=\alph*)]
  \item Calcule la masse $m=m_2-m_1$ du précipité avec le bon nombre de décimales.
  \item Compare l'incertitude relative sur $m_1$ et celle sur $m$. Que constate-t-on ?
\end{enumerate}
\end{exo}

\begin{exo}[la même balance, deux échantillons]
Une balance est précise au centigramme ($\Delta m=\qty{0,01}{\gram}$). On pèse un échantillon
d'environ $\qty{2}{\gram}$, puis un autre d'environ $\qty{200}{\gram}$.
\begin{enumerate}[label=\alph*)]
  \item Calcule l'incertitude relative dans chaque cas.
  \item Pour quelle pesée la balance est-elle suffisante si l'on vise une incertitude relative
        $\le 0,1\%$ ? Conclus.
\end{enumerate}
\end{exo}

\end{document}
